Une corde de \qty{\v{L}}{m}, fixée à ses deux extrémités et de masse linéique \qty{\v{mu}}{g/m}, est mise en vibration par un vibreur de fréquence réglable. On observe deux résonances successives aux fréquences \qty{\fpeval{\v{n}*\v{f1}}}{Hz} et \qty{\fpeval{(\v{n}+1)*\v{f1}}}{Hz}, et aucune résonance entre les deux.

\begin{enumerate}[label=\alph*)]
    \item Déterminez la fréquence fondamentale de la corde.
    \item À quels modes ($n$) correspondent les deux résonances observées ?
    \item Déterminez la vitesse des ondes sur la corde.
    \item Déterminez la tension dans la corde.
\end{enumerate}
